\origpage{78}{ROBERTSON AND SEYMOUR}
\hspace*{12pt}(ii) there exist curves $\phi, \phi'$ tracing $W, W'$, respectively, such that $\phi$ is homotopic in $Z$ to $\phi'$.

A walk $W$ is \textit{null-homotopic in} $Z$ if it is homotopic in $Z$ to a walk of length zero (and hence is closed).

In all the forms of homotopy mentioned above, we shall usually abbreviate ``homotopic in $\Sigma$'' to ``homotopic.''

\sectionhead{4. Slopes}

Let $\Sigma$ be a surface, and let $\Gamma$ be a painting or drawing in $\Sigma$. A subset $Z \subseteq \Sigma$ is \textit{$\Gamma$-normal} if $Z \cap U(\Gamma) \subseteq V(\Gamma)$. If $\Sigma$ is connected, we say that $\Gamma$ is \textit{$\theta$-representative} if every $\Gamma$-normal $O$-arc $F$ with $|F \cap V(\Gamma)| < \theta$ is null-homotopic. If $\Delta \subseteq \Sigma$ is closed and $bd(\Delta)$ is $\Gamma$-normal then $(U(\Gamma) \cap \Delta, V(\Gamma) \cap \Delta)$ is a painting or drawing (respectively) which we denote by $\Gamma \cap \Delta$. A tangle $\mathscr{T}$ of order $\theta$ in a painting or drawing $\Gamma$ is said to be \textit{respectful} if $\Sigma$ is connected and for every $\Gamma$-normal $O$-arc $F$ in $\Sigma$ with $|F \cap V(\Gamma)| < \theta$, there is a closed disc $\Delta \subseteq \Sigma$ with $bd(\Delta) = F$ such that
\begin{equation*}
(\Gamma \cap \Delta, \Gamma \cap \overline{\Sigma - \Delta}) \in \mathscr{T}.
\end{equation*}
If there is a respectful tangle of order $\theta$ then clearly $\Gamma$ is $\theta$-representative, and our first main result is the converse for non-spherical $\Sigma$, the following.

\stmt
(4.1) \textit{Let $\Sigma$ be a connected surface, not a sphere, let $\theta \ge 1$, and let $\Gamma$ be a $\theta$-representative $2$-cell painting or drawing in $\Sigma$. Then there is a unique respectful tangle in $\Gamma$ of order $\theta$.}

Statement (4.1) will be proved in Section 6. To prove (4.1) we use the concept of a ``slope'' in a drawing $\Gamma$ in a connected surface $\Sigma$, defined as follows. Let $\theta \ge 1$ be half an integer. A \textit{slope} in $\Gamma$ of order $\theta$ is a function $\text{ins}$ which assigns to every circuit $C$ of $\Gamma$ of length $< 2\theta$ a closed disc $\text{ins}(C) \subseteq \Sigma$, bounded by $U(C)$, such that

\hspace*{12pt}(i) if $C_1, C_2$ are circuits of $\Gamma$, both of length $< 2\theta$, and $U(C_1) \subseteq \text{ins}(C_2)$, then $\text{ins}(C_1) \subseteq \text{ins}(C_2)$

\hspace*{12pt}(ii) if $P_0, P_1, P_2$ are three paths of $\Gamma$ of length $> 0$, with the same ends, mutually internally disjoint, and the circuits $P_1 \cup P_2$, $P_2 \cup P_0$, $P_0 \cup P_1$ all have length $< 2\theta$, then one of
\begin{equation*}
\text{ins}(P_1 \cup P_2), \quad \text{ins}(P_2 \cup P_0), \quad \text{ins}(P_0 \cup P_1)
\end{equation*}
includes the other two.

(Paths and circuits have no ``repeated'' vertices, and the length of a path or
